\section{Introduction: the project and what has been learned}
\label{sec:introduction}
The purpose of this project is to explain an arithmetic positivity statement by constructing a system whose positivity is already understood. The arithmetic object is the Weil quadratic form associated with the Riemann zeta function. Its nonnegativity on every smooth compactly supported complex input is equivalent to the Riemann hypothesis. The desired argument would identify this form with an ordinary positive bulk energy or Hilbert-space norm. Numerical tests could then check the identity; they would not be the source of its sign.

This manuscript develops one route toward such a construction. It studies supersymmetric ground states, their parameter-dependent metrics, and interacting theories that retain prescribed gamma and prime factors. It produces exact local-factor identities and exclusions of specified repairs. It does not produce the full Weil norm identity. Understanding both the successes and their limits is essential: a positive quantum Hamiltonian can encode an arithmetic factor without making the required arithmetic quadratic form positive.

\subsection{The arithmetic quantity that the theory must reproduce}
Let $f$ be a smooth complex function supported in $I_L=(-L/2,L/2)$, and let $F$ be its zero extension to $\R$. Use the Fourier convention
\[
\widehat F(\tau)=\int_{\R}F(x)e^{-i\tau x}\,dx.
\]
In the normalization used here, the target has the decomposition
\begin{equation}\label{intro:target}
\begin{split}
Q_L[f]={}&K[F]+w_0\norm f^2+2|C(f)|^2-2|S(f)|^2\\
&-2\sum_{\substack{p\ {
m prime},\ m\ge1\\m\log p<L}}
(\log p)p^{-m/2}\re\ip F{U_{m\log p}F}.
\end{split}
\end{equation}
Here $U_dF(x)=F(x-d)$ compares the profile with a translated copy, and
\[
C(f)=\int_{I_L}f(x)\cosh(x/2)\,dx,\qquad
S(f)=\int_{I_L}f(x)\sinh(x/2)\,dx.
\]
The two squared global amplitudes constitute the signed pole contribution. The gamma kinetic energy and its local normalization are
\begin{align}
K[F]&=\frac1{2\pi}\int_{\R}B(\tau^2)|\widehat F(\tau)|^2\,d\tau,
& w_0&=\psi(1/4)-\log\pi<0,\label{intro:kinetic}\\
B(s)&=\re\psi(1/4+i\sqrt{s}/2)-\psi(1/4),\qquad s\ge0,\nonumber
\end{align}
where $\psi=\Gamma'/\Gamma$ is the digamma function. The distances $m\log p$ and weights $(\log p)p^{-m/2}$ are fixed arithmetic data. A finite support length activates only finitely many such distances.

The gamma kinetic term is nonnegative, but the contact constant, translated correlations and pole contribution cannot be discarded. The project requires the complete sum in \eqref{intro:target} for all inputs and arbitrarily large $L$, rather than a chosen finite collection of functions or a fixed list of primes at every length. The Weil criterion and modern operator formulations provide the mathematical setting \cite{Suzuki}. Section~\ref{sec:arithmetic} supplies the form domain and the identities used below; no separate background document is needed to read this manuscript.

\subsection{What a positive bulk or quantum realization would mean}
In a classical auxiliary-field construction, one fixes the source $f$, minimizes a nonnegative bulk energy over responding fields $u$, and seeks
\begin{equation}\label{intro:bulk-goal}
Q_L[f]=\inf_u\mathcal E_L[f,u],\qquad \mathcal E_L[f,u]\ge0.
\end{equation}
The bulk may consist of fields on an extra coordinate, a collection of channels, or other auxiliary degrees of freedom. Its defining action and source map must be specified independently of the unknown positivity of $Q_L$.

The operator version seeks a positive physical Hilbert space and a linear preparation map $\mathcal A_L$ with
\begin{equation}\label{intro:norm-goal}
Q_L[f]=\norm{\mathcal A_Lf}_{\mathscr H}^2.
\end{equation}
Equivalently, the associated form operator would satisfy $W_L=\mathcal A_L^*\mathcal A_L$ in the appropriate form sense. This replaces separate tests of individual functions by one operator identity. The requirement that the identity hold on the full input domain remains. Defining $\mathcal A_L$ through an assumed positive square root of $W_L$ would simply assume the desired result.

Supersymmetry offers a way to organize such a physical construction. A nilpotent charge $Q$ with its actual Hilbert adjoint gives
\[
H=QQ^\dagger+Q^\dagger Q\ge0.
\]
Its zero-energy states are harmonic representatives: states annihilated by both charges. Their ordinary norms can be nonzero even though their energies vanish. Thus the proposed observable in a supersymmetric ground-state construction is a prepared norm or boundary pairing, not merely the ground-state energy. A positive Hamiltonian alone does not establish \eqref{intro:norm-goal}. Quantizing an auxiliary system is useful only if the quantum state space, interactions or gluing law supplies an identity absent from the classical minimization.

\subsection{The starting success and the remaining discrepancy}
The gamma kinetic term already has an exact positive auxiliary-field realization. The identity
\[
B(s)=\sum_{k\ge0}\frac2{a_k}\frac{s}{s+a_k^2},\qquad a_k=2k+\tfrac12,
\]
expresses it as a positive tower of resolvent responses. Each response comes from minimizing a squared mismatch and a positive gradient energy. This is a substantial first success: an entire continuous operator, on its logarithmic form domain, is represented by a positive field energy. It does not include the contact or poles.

A second starting model is a positive supersymmetric Morse system tensored with discrete prime sectors. For a finite prime set $\mathcal S$ it has a state $\Psi_\sigma$ and a self-adjoint observable $X$ satisfying
\begin{equation}\label{intro:amplitude}
\ip{\Psi_\sigma}{e^{i\tau X}\Psi_\sigma}
=\Lambda_{\mathcal S}(\sigma+i\tau),\qquad
\Lambda_{\mathcal S}(z)=\pi^{-z/2}\Gamma(z/2)
\prod_{p\in\mathcal S}(1-p^{-z})^{-1}.
\end{equation}
The formula is verified by an elementary integral and geometric sums. Its normalized logarithmic derivative yields a positive gamma-plus-prime jump energy $J_{\mathcal S,L}$. The complete target, however, is
\begin{equation}\label{intro:missing}
W_L=J_{\mathcal S,L}-\kappa_{\mathcal S}I+P_L,
\qquad \kappa_{\mathcal S}>0,
\end{equation}
when $\mathcal S$ includes all primes active on the interval. The missing negative contact and signed pole operator are therefore precise, rather than unspecified corrections. A positive characteristic function or jump process does not by itself account for them.

\subsection{Why enlarge the theory and admit interactions?}
The original arithmetic state family is too restrictive for the desired extended supersymmetry. Its chosen complex extension has a nonflat scalar metric, whereas an ordinary one-vacuum chiral $tt^*$ system has zero curvature. Here $tt^*$ denotes the compatibility equations relating a physical vacuum metric, its projected connection, and operators induced by variations of a holomorphic superpotential. More vacua permit matrix commutators and hence additional geometry. The underlying fields and dynamics must actually support those equations; changing the name of a two-charge model does not create four physical supercharges.

We therefore allow new bosons, fermions, boundary data and genuine interactions. The models studied below are supersymmetric quantum mechanics with complex coordinates and a holomorphic superpotential, often called Landau--Ginzburg models. Their fermions are represented by differential forms, and their Hamiltonian positivity follows from physical adjoints. A polynomial superpotential can create several isolated vacua and a nontrivial matrix of ground-state overlaps. A periodic complex coordinate introduces an additional global issue: some parameter insertions become multivalued, so ordinary chiral $tt^*$ equations cannot automatically be applied without the periodic or flavor data \cite{CGV}.

Gaussianity is consequently a modeling choice, not the boundary between possible and impossible arithmetic constructions. Quadratic models are valuable controls because they can be solved exactly. An interacting theory is more useful only when its interaction changes the pairing or constraints relevant to \eqref{intro:norm-goal}. The calculations here give a concrete example of that distinction.

\subsection{Overview of the results}
First, replacing the real Morse coordinate by a complex cylinder produces an actual four-supercharge theory. It has a normalizable one-form vacuum for nonzero cylinder parameter. The gamma function survives as a holomorphic integral, but this integral is not identified with the physical norm of that vacuum. Two added quadratic complex prime fields give both an Euler period and a physical metric proportional to $|1-q|^{-2}$ in a specified holomorphic frame. The latter has zero local curvature, so its interpretation depends on the physical choice of source.

Second, a cross-coupled quartic prime theory has nine normalizable middle-degree vacua. At fixed nonzero quartic coupling, an exact differential equation excludes the Euler identity period on every flat rapid-decay cycle obeying the stated boundary conditions. This is an analytical exclusion of a class of theories and observables, not a failure to fit a numerical curve.

Third, scaling the quartic coupling with the square of the prime mass restores the Euler factor exactly. With $M=1-q$ and a fixed interaction shape $c>0$, the identity period and physical identity-class metric have the forms
\begin{equation}\label{intro:interacting-success}
\mathcal P_c(M)=\frac{C_c}{M},\qquad
h_c(M,\bar M)=\frac{H_c}{|M|^2},\qquad C_c,H_c>0.
\end{equation}
This shows that an interacting positive theory can carry the exact Euler dependence. The constants describe different quantities and are not known to satisfy $H_c=|C_c|^2$. Moreover, the mass dependence is generated by a field dilation: the mass deformation vanishes in the Jacobi ring, the algebra of holomorphic observables modulo the superpotential gradients. It therefore supplies no new ordinary chiral-curvature constraint along that parameter.

Finally, several apparently simple completions are excluded under explicit hypotheses. Finite ambient-holomorphic vacuum enlargement conflicts with the trace of the $tt^*$ curvature equation. Scalar reversible ground transforms and scalar phase changes have the wrong off-diagonal signs for the full target. A positive square representing the pole correlations necessarily adds a diagonal energy, which must also be accounted for. Curvature equations alone leave a normalization freedom that changes the contact term. These restrictions still allow matrix or fermionic interference, nonlocal source maps and physically specified relative boundary pairings.

The resulting next problem is narrower than a search for an arbitrary complicated quantum theory. Starting with the exact interacting control \eqref{intro:interacting-success}, one should change the interaction shape so that its parameter operator is nonzero in the Jacobi ring, and identify the physical boundary pairing before attempting a large overlap calculation. A source-normalization-invariant comparison between its physical metric and squared period is one concrete diagnostic. The decisive advance would be a pairing identity that supplies the contact and poles together with the full input space. No such identity, all-support positivity theorem or Riemann-hypothesis proof is claimed here.

\subsection{Organization and status}
Section~\ref{sec:arithmetic} derives the arithmetic starting identities. The subsequent sections construct the enlarged theory, analyze the interacting Euler sector, and develop the metric and completion restrictions. All numbered results specify their model and hypotheses. The mathematical and physical frameworks are attributed to external primary references; the particular algebraic tests and model comparisons are derived explicitly. This is an AI-assisted research draft. Internal cross-review and exact algebra checks do not replace independent specialist review or establish a priority claim.
